\documentclass[11pt]{article}
\usepackage{mathunal-base}
\mathunalfuente{serif}

\renewcommand{\examtitulo}{Cálculo en Varias Variables --- Tercer Parcial}
\renewcommand{\examfecha}{Semestre 01-2018, 28 de mayo de 2018}

\begin{document}

\examheader

\vspace{4pt}
\noindent Nombre: \dotfill\ Cédula: \dotfill

\vspace{4pt}
\noindent Grupo: \dotfill\ Salón: \dotfill

\vspace{4pt}
\noindent Puntaje. Solo para uso oficial:\quad 1--4 (/28) \dotfill;\quad 5 (/24) \dotfill;\quad 6 (/24) \dotfill;\quad 7 (/24) \dotfill.\quad Total: \dotfill\quad Nota: \dotfill

\vspace{6pt}
\noindent\textbf{Instrucciones.} Antes de empezar verifique que su examen esté completo y que sus datos sean correctos. Por favor verifique que su celular esté apagado. No está permitido el uso de calculadora ni de otros aparatos electrónicos. Solucione las preguntas en los espacios en blanco asignados a cada una de ellas. No está permitido sacar hojas en blanco ni ningún tipo de apuntes durante el examen. Duración: 1 hora y 50 minutos.

\vspace{8pt}
\noindent\textit{En las preguntas 1 a 4 responda en el espacio en blanco. En estas preguntas la calificación tendrá en cuenta sólo la respuesta.}

\begin{enumerate}
\item{} (7\,\%) Sea $C$ la frontera del triángulo con vértices $(-1,0)$, $(0,1)$ y $(1,0)$, orientada en sentido antihorario. El valor de la integral
\[
\int_C \left(\operatorname{sen}(e^x) - 3y\right)dx + \left(\tan^{-1}\!\left(\frac{y^2}{y^4+1}\right) + 10x\right)dy
\]
es \underline{\phantom{XXXXXX}}.

\item{} (7\,\%) Un alambre delgado está doblado en forma de una semicircunferencia $x^2 + y^2 = 9$, $x \geq 0$. Si la densidad lineal es $\rho(x,y) = x$ entonces la masa del alambre es \underline{\phantom{XXXXXX}}.

\item{} (7\,\%) El flujo hacia afuera del campo vectorial $F(x,y,z) = (xy^2,\ x^2y,\ y)$ a través de la superficie $S$ que es la frontera de la región encerrada por el cilindro $x^2 + y^2 = 1$ y los planos $z = -1$ y $z = 1$ es \underline{\phantom{XXXXXX}}.

\item{} (7\,\%) Unas ecuaciones paramétricas para la recta tangente a la curva con ecuaciones paramétricas
\[
x = 2\cos t,\quad y = \operatorname{sen} t,\quad z = t
\]
en el punto $(0, 1, \pi/2)$ son
\[
x = \underline{\phantom{XXXX}},\qquad y = \underline{\phantom{XXXX}},\qquad z = \underline{\phantom{XXXX}}.
\]
\end{enumerate}

\vspace{4pt}
\noindent\textsc{Preguntas con procedimiento.} En las preguntas 5 a 7 responda mostrando \textit{detalladamente} el procedimiento utilizado.

\begin{enumerate}[start=5]
\item{} (24\,\%) Calcule el área superficial del paraboloide $z = 3 - x^2 - y^2$ que está por encima de $z = 1$.

\item{} (24\,\%) Sea $S$ la porción del plano $x + y + z = 1$, en el primer octante, orientada con vector normal unitario apuntando hacia arriba, y sea $C$ la frontera de $S$. Calcule el trabajo realizado por el campo de fuerzas
\[
F(x,y,z) = \left(z - \cos(e^x),\ x + e^{y^4},\ y + \sqrt{z^2+1}\right),
\]
para desplazar una partícula a lo largo de $C$.

\item{} (24\,\%) Calcule el flujo hacia afuera de la esfera $x^2 + y^2 + z^2 = z$ del campo vectorial
\[
F(x,y,z) = \left(\frac{x^3}{3} - \cos(e^y),\ yz^2 + e^{x^4},\ zy^2 + \tan^{-1}(y^2)\right).
\]
\end{enumerate}

\examfooter

\end{document}
